\section*{An annealed upper-set envelope for critical restrictions}

All logarithms are base two.  Let $f:\{-1,1\}^n\to\{-1,1\}$ and put
\[
 p(T)=\widehat f(T)^2,\qquad
 \pi(A)=\sum_{T\supseteq A}p(T).
\]
Let $J\subseteq[n]$ retain every coordinate independently with probability
$t$, and let $X$ be a uniform assignment of $J^c$.  Write
$g=f_{J,X}$.  For $A\subseteq J$ define
\[
 q_{J,X}(A)=\widehat g(A)^2,\qquad
 \Pi_{J,X}(A)=\sum_{C:\,A\subseteq C\subseteq J}q_{J,X}(C),
\]
and the phase-free marginal
\[
 m_J(A)=\sum_{T:\,T\cap J=A}p(T).
\]

\begin{theorem}[annealed upper-set envelope]
For every Boolean $f$ and every $t\in[0,1]$,
\[
 \boxed{\quad
 \mathbb E_{J,X}U(f_{J,X})\le F_t(p):=
 \mathbb E_J\sum_{A\subseteq J}m_J(A)
       \log\frac{\pi(A)}{m_J(A)}.\quad}                 \tag{1}
\]
Terms with zero mass vanish.  This is an all-dimensional statement; it does
not assume criticality or a lower bound on any spectral marginal.
\end{theorem}

\begin{proof}
The restriction Fourier formula gives
\[
 \widehat g(A)=\sum_{B\subseteq J^c}\widehat f(A\cup B)\chi_B(X),
\]
so orthogonality in $X$ yields
\[
 \mathbb E_Xq_{J,X}(A)=m_J(A).                           \tag{2}
\]
Consequently
\[
 \mathbb E_X\Pi_{J,X}(A)
 =\sum_{C:\,A\subseteq C\subseteq J}m_J(C)
 =\sum_{T\supseteq A}p(T)=\pi(A).                       \tag{3}
\]
On $0\le x\le y$, the continuous extension of
$\phi(x,y)=x\log(y/x)$ is jointly concave.  Indeed, after multiplying its
Hessian by $\ln2$, one obtains
\[
 \begin{pmatrix}-1/x&1/y\\1/y&-x/y^2\end{pmatrix},
\]
which is negative semidefinite (its determinant is zero).  Jensen's inequality,
(2), and (3) therefore give
\[
 \mathbb E_X\phi(q_{J,X}(A),\Pi_{J,X}(A))
 \le\phi(m_J(A),\pi(A)).
\]
Summing over $A\subseteq J$ and then averaging $J$ proves (1). \qed
\end{proof}

\subsection*{Exact coordinate charge}

Fix once and for all an order on $[n]$.  Let $Y$ be an independent spectral
sample with law $p$.  For $A\subseteq J$ and $i\in J\setminus A$, put
\[
 Z_{<i}=\{j\in J\setminus A:j<i\}
\]
and
\[
 d_{J,A,i}=m_J(A)\log
 \frac{\Pr[Y\supseteq A,\;Y\cap Z_{<i}=\varnothing]}
      {\Pr[Y\supseteq A,\;Y\cap(Z_{<i}\cup\{i\})=\varnothing]}.
                                                               \tag{4}
\]
Every denominator is positive whenever $m_J(A)>0$.  Set
\[
 C_i(t)=\mathbb E_J\sum_{A\subseteq J\setminus\{i\}}d_{J,A,i}.
\]
The probabilities in (4) telescope from $\pi(A)$ to $m_J(A)$ as the
coordinates of $J\setminus A$ are exposed.  Hence
\[
 F_t(p)=\sum_{i=1}^n C_i(t),\qquad C_i(t)\ge0.             \tag{5}
\]
Writing $p_i=\Pr[i\in Y]=\operatorname{Inf}_i(f)$, Parseval gives
$I(f)=\sum_i p_i$.  If $p_i=0$, then $C_i(t)=0$.  Thus, for every nonconstant
$f$, some physical coordinate satisfies
\[
 \boxed{\qquad \frac{C_i(t)}{p_i}\ge\frac{F_t(p)}{I(f)}.\qquad}          \tag{6}
\]
The charge has the following probabilistic meaning: sample $J$, sample the
observed intersection $A=Y\cap J$, and charge each retained absent coordinate
the conditional surprisal of its absence after the earlier absent coordinates
and all observed present coordinates have been revealed.

\subsection*{Consequence for finite-dimensional ratio maximizers}

Let $f$ attain
\[
 c_n=\max_{h\text{ nonconstant on }n\text{ bits}}H(h)/I(h),
 \qquad R=c_n.
\]
The source-complete A052 random-restriction persistence theorem states
\[
 \mathbb E U(g)\ge
 \left[\frac{R-2}{2}t-\frac{1-t}{\ln2}\right]I(f)
 +\mathbb E\operatorname{Var}(g)+\frac12\mathbb E b(g),                 \tag{7}
\]
where
\[
 b(g)=\widehat g(\varnothing)^2
       \log\frac1{\widehat g(\varnothing)^2}.
\]
Combining (1), (6), and (7), at $t=1/2$ there is an $i$ with $p_i>0$ such
that
\[
 \boxed{
 \frac{C_i(1/2)}{p_i}\ge \frac{F_{1/2}(p)}{I(f)}
 \ge \frac{R-2}{4}-\frac1{2\ln2}
 +\frac{\mathbb E\operatorname{Var}(g)+\frac12\mathbb E b(g)}{I(f)}.}  \tag{8}
\]
In particular, any hypothetical sequence of finite-dimensional maximizers
with $R\to\infty$ must contain a physical coordinate whose normalized
conditional-absence surprisal diverges at least as $R/4-O(1)$.  By the prior
all-block critical theorem, once $R>4$ every such maximizer is already
spectrally centered: $p_i<1/2$ for every nontrivial coordinate (more precisely,
$\Psi(p_i)/p_i\ge R$ in that theorem's notation).  Thus the witness in (8)
cannot be dismissed as a high-influence or spectrally determined coordinate.

\subsection*{Scope and remaining payment}

Equation (8) is a new necessary structural condition on maximizers and an
exact phase-free upper envelope for the restriction frontier.  It is not an
upper bound $U\le KI$ and therefore does not prove FEI.  The precise unpaid
step is now static: prove, for a ratio-critical Boolean spectrum, that
$\max_i C_i(1/2)/p_i\le K$ with universal $K$ (or at least bound the
maximizing witness selected in (6)), or construct a genuine Boolean critical
family for which this maximum diverges.
Arbitrary probability distributions are not admissible substitutes for such a
Boolean family.
